\documentclass[10pt,a4paper]{amsart}
\usepackage[margin=19mm]{geometry}
\usepackage{amsmath,amssymb,mathtools}
\usepackage[scr=rsfs,cal=euler]{mathalfa}
\usepackage{xfrac}
\usepackage[unicode,hidelinks,bookmarks=false]{hyperref}
\newcommand{\ie}{i.e.\ }
\newcommand{\cf}{cf.\ }
\newcommand{\resp}{respectively\ }
\newcommand{\Subsection}{\subsection}
\providecommand{\nobreakdash}{\nobreak\hbox{-}}
\setlength{\emergencystretch}{2em}
\multlinegap0pt
\usepackage{bm}
\usepackage{bbm}
\usepackage{dsfont}
\usepackage{xspace}
\usepackage{cancel}
\usepackage{mathtools}
\usepackage{amsthm}
\makeatletter
\@ifundefined{assumption}{\newtheorem{assumption}{Assumption}}{}
\@ifundefined{definition}{\newtheorem{definition}{Definition}}{}
\@ifundefined{theorem}{\newtheorem{theorem}{Theorem}}{}
\makeatother
\newcommand{\K}{\mathcal{K}}
\newcommand{\R}{\mathbb{R}}
\newcommand{\Supp}{\mathrm{Supp}}
\newcommand{\X}{\mathcal{X}}
\newcommand{\cB}{\mathcal{B}}
\newcommand{\tgt}{\mathrm{tgt}}
\title[Symplectic Inductive Bias for Data-Driven Target Reachabilit]{Symplectic Inductive Bias for Data-Driven Target Reachability in Hamiltonian Systems}
\author{Zhuo Ouyang, Jixian Liu,, Enrique Mallada}
\date{Source: arXiv:2604.17213v1 (CC BY 4.0)}
\begin{document}
\maketitle

\noindent\emph{Source and licence.} Symplectic Inductive Bias for Data-Driven Target Reachability in Hamiltonian Systems. Version arXiv:2604.17213v1 (CC BY 4.0), distributed under the Creative Commons Attribution 4.0 International licence, as recorded in the arXiv licence metadata for this exact version and shown on the abstract page https://arxiv.org/abs/2604.17213v1. Full rights notice: licenses/arxiv-2604.17213-CC-BY-4.0.txt.\par\smallskip

\noindent\textit{Editorial scope.} Counted unit: the Theorem whose source label is \texttt{thm\_existence\_chain\_policy} in the article (source0.tex lines 605--620, source label \texttt{thm\_existence\_chain\_policy}), delivered together with the article's own complete original proof of it (source0.tex lines 622--789, one \texttt{proof} environment of 168 lines). No mathematics is written, repaired or re-derived: statement and proof are the article's own text line for line. Required context retained uncounted: eq:Hamil\_sys (lines 173--180); ass:bounded\_gradient (lines 185--188); ass\_reachability\_of\_the\_target\_set (lines 399--405); ass\_lowest\_speed (lines 427--433); thm\_general\_reachability (lines 438--468); condition1 (lines 444--462); condition2 (lines 447--450); condition3 (lines 447--450); ass\_dim\_ergodic\_components (lines 579--585); ass\_strong\_convexity (lines 587--593). Every definition, display and label of those lines is retained unchanged. Omitted from this package and not redistributed: 1--172, 181--184, 189--398, 406--426, 434--437, 451--578, 586--586, 594--604, 621--621, 790--1039. Nothing inside a retained excerpt is omitted or altered: every retained body is delivered exactly as the article's lines are written, so no omission receipt of any kind accompanies this package. Citations: the retained excerpts carry no citation command at all, so no bibliography entry of the article is transcribed and no reference is redistributed. Reference adaptation: none was needed and none was made; every reference inside the delivered unit resolves inside it, so no unresolved reference or citation is produced. Printed slips kept literal: no printed word, symbol, hypothesis or number of the counted statement or of its proof was altered. Layout and class: the article's own class is \texttt{ieeeconf} with packages including inputenc, algorithm, algpseudocode, amsmath, amssymb, amsthm, graphicx, subcaption, booktabs, placeins, hyperref, multirow, caption, todonotes; the wrapper is the lane's pinned \texttt{amsart} editorial wrapper at 10pt on A4, which restates the theorem-like environments the retained text needs and adds the article's own macro definitions that the retained text needs (\texttt{\textbackslash K}, \texttt{\textbackslash R}, \texttt{\textbackslash Supp}, \texttt{\textbackslash X}, \texttt{\textbackslash cB}, \texttt{\textbackslash tgt}), with extra wrapper packages (bm, bbm, dsfont, xspace, cancel, mathtools). The delivered numbering is the unit's own. Assets: no retained span contains a figure environment, a graphics-inclusion command, a PDF-page inclusion, a table or a TikZ picture; no image file is copied into this package, the package declares no asset dependency and no image file is redistributed with it. Licence: version arXiv:2604.17213v1 of this article is distributed under the Creative Commons Attribution 4.0 International licence, as recorded in the arXiv licence metadata for this exact version and shown on the abstract page https://arxiv.org/abs/2604.17213v1. A full rights notice is delivered at licenses/arxiv-2604.17213-CC-BY-4.0.txt. Characters: every retained excerpt is pure ASCII, so the delivered engine, XeLaTeX with its default Latin Modern text font, reports no missing character.
\begin{definition}[Hamiltonian System]
A Hamiltonian system without energy dissipation is a dynamical system of the form
\begin{align}
\label{eq:Hamil_sys}
    \dot{x} = f(x,u) = J(x)\nabla H(x) + G(x)u,
\end{align}
where $x \in \X \subset \R^n$ is the state defined on a compact set $\X$, $H:\R^n \to \R$ is the Hamiltonian function, $J(x)\in\R^{n\times n}$ is skew-symmetric, and $u\in U\subset \R^m$ is the input of the system defined on a compact U.
\end{definition}

\begin{assumption}[Bounded Hamiltonian Gradient]
\label{ass:bounded_gradient}
The Hamiltonian function $H$ is continuously differentiable on $\X$ and there exists an upper bound $L_H > 0$ such that $\forall x\in\X, \|\nabla H(x)\|\le L_H$.
\end{assumption}

\begin{assumption}[Reachability of $H_{\tgt}^\epsilon$]
\label{ass_reachability_of_the_target_set}
For all $x \in \mathcal{X} \setminus H_{\tgt}^\epsilon$, there exist $T > 0$ and $u \in \mathcal{U}^{(0,T]}$ such that
\[
\phi(T,x,u) \in H_{\tgt}^\epsilon.
\]
\end{assumption}

\begin{assumption}[Uniform Positive Energy Decrease Rate]
\label{ass_lowest_speed}
There exists $\epsilon>0$ such that
\[
\underline{v_\epsilon} := \inf_{x\in \mathcal{X}\setminus H_{\tgt}^\epsilon} v_\epsilon(x) > 0.
\]
\end{assumption}

\begin{theorem}[Target Reachability]
\label{thm_general_reachability}
Consider system~\eqref{eq:Hamil_sys} under Assumptions~\ref{ass:bounded_gradient}--\ref{ass_lowest_speed}. 
Let $\mathcal{K} = \{(x_i, r_i, u_i)\}_{i=1}^N$ be a finite assignment set. 
Assume that $\mathcal{K}$ satisfies the following:

\begin{enumerate}
\item \textbf{Local energy decrease:}\label{condition1}
For each $(x_i,r_i,u_i)\in \mathcal{K}$,
\begin{align*}
\Delta H(x_i(\tau_i)) \!+\! v_0 \tau_i \!+\! L_H r_i e^{L \tau_i}
\!\le\! \Delta H(x_i) \!-\! L_H r_i,
\end{align*}
where $x_i(\tau_i):=\phi(\tau_i,x_i,u_i)$, $\tau_{\min}=\min_i\tau_i,\ \text{and }v_0 \textgreater 0$.
\item \textbf{Energy coverage:}\label{condition2}
Let $c := \sup_{x\in S_0} \Delta H(x)$. Then
\[
\Delta H(x)\le c \;\implies\; H(x)\in H(\Supp(\mathcal{K})).
\]
\item \textbf{Ergodic coverage:}\label{condition3}
For all $E \in H(\Supp(\mathcal{K}))$ and all ergodic components $K_\alpha^E \subseteq \Sigma_E$,
\[
\mathrm{int}(\Supp(\mathcal{K})) \cap K_\alpha^E \neq \emptyset.
\]
\end{enumerate}

Then, the chain policy $\pi_{\mathcal{K}}$ ensures that for almost every $x_0 \in S_0$, there exists $t<\infty$ such that
\[
\phi(t,x_0,\pi_{\mathcal{K}}) \in S_{\tgt}.
\]
\end{theorem}

\begin{assumption}[Ergodicity of Energy Layers]
\label{ass_dim_ergodic_components}
For every $E \in H(\Supp(\mathcal{K}))$, the energy layer $\Sigma_E$ is ergodic, i.e.,
\[
\forall E \in H(\Supp(\mathcal{K})),\quad K_\alpha^E = \Sigma_E.
\]
\end{assumption}

\begin{assumption}[Strong Convexity of the Hamiltonian]
\label{ass_strong_convexity}
The Hamiltonian function $H$ is strongly convex on $\mathcal{X}$, i.e., there exists $\mu_H>0$ such that for all $x,y\in\mathcal{X}$,
\[
H(y)\ge H(x)+\nabla H(x)^\top (y-x)+\frac{\mu_H}{2}\|y-x\|^2.
\]
\end{assumption}

\begin{theorem}[Existence of Chain Policy]
\label{thm_existence_chain_policy}
Consider system~\eqref{eq:Hamil_sys} satisfying Assumptions~\ref{ass_reachability_of_the_target_set}--Assumption~\ref{ass_strong_convexity}. Let $v_0 \in (0,\underline{v_\epsilon})$.
Then there exists a nonparametric chain policy $\pi_{\mathcal{K}}$, constructed from a finite assignment set $\mathcal{K}=\{(x_i,r_i,u_i)\}_{i=1}^{N}$, where
\begin{align*}
r_i=\frac{(v_\epsilon(x_i)-v_0)T_\epsilon^\star(x_i)}{L_H(1+e^{LT_\epsilon^\star(x_i)})}>0,
\end{align*}
with $u_i:(0,T_\epsilon^\star(x_i)]\to {U}$ being the optimal control for reaching $\partial H_{\tgt}^\epsilon$ from $x_i$, $T_\epsilon^\star(x_i)$ denoting its hitting time, such that the following holds:
\begin{enumerate}
\item \textbf{Reachability:} For almost every $x_0\in S_0$, there exists $t<\infty$ such that $\phi(t,x_0,\pi_{\mathcal{K}})\in S_{\tgt}$.
\item \textbf{Sample complexity:} Let $H(\Supp(\mathcal{K}))=[H_1,H_2]$. Then
\[
N\leq (H_2-H_1)\frac{16L_H^2}{\mu_H (1-\frac{v_0}{\underline{v_\epsilon}})^2}\frac{\exp(\frac{2L(L_HD_{\mathcal{X}}+\epsilon)}{\underline{v_\epsilon}})}{\epsilon^2}.
\]
\end{enumerate}
\end{theorem}

\begin{proof}
We construct a canonical assignment set and then verify the three conditions of Theorem~\ref{thm_general_reachability}.

\medskip
\noindent\textbf{Step 1: Canonical local construction.}
Let 
\[
\X_c:=\{x\in \X:\Delta H(x)\le c\}\setminus H_{\tgt}^\epsilon.
\]
For every $x \in \X_c$, by Assumptions~\ref{ass_reachability_of_the_target_set}
and~\ref{ass_lowest_speed}, choose an optimal control $u_x:(0,T_\epsilon^\star(x)] \to U$, choose an optimal control $u_x:(0, T_\epsilon^\star(x)] \to U$ such that 
\[
\phi(T_\epsilon^\star(x),x,u_x) \in \partial H_{\tgt}^\epsilon, \,
v_\epsilon(x) = \frac{\Delta H(x) + \epsilon}{T_\epsilon^\star(x)}
\]
Define the certified radius
\[
r(x)
:=
\frac{(v_\epsilon(x)-v_0) T_\epsilon^\star(x)}
{L_H\bigl(1+e^{LT_\epsilon^\star(x)}\bigr)}.
\]
Since $v_0<\underline v_\epsilon\le v_\epsilon(x)$, we have $r(x)>0$.
Moreover,
\begin{align*}
 & \Delta H(\phi(T_\epsilon^\star(x),x,u_x))
+v_0T_\epsilon^\star(x)
+L_Hr(x)e^{LT_\epsilon^\star(x)}\\
= &
\Delta H(x)-L_Hr(x),
\end{align*}
so each triple $(x,r(x),u_x)$ satisfies Condition~\ref{condition1} of
Theorem~\ref{thm_general_reachability}.

\noindent\textbf{Step 2: Uniform lower bound on the certified radii.}
According to Assumption~\ref{ass_lowest_speed}: $v_\epsilon(x)\geq \underline{v_\epsilon}>0,
v_\epsilon(x)=\frac{\Delta H(x_i)+\epsilon}{T_\epsilon^\star(x)},$
we have $T_\epsilon^\star(x) \leq \frac{\Delta H(x)+\epsilon}{\underline{v_\epsilon}}$.
Now recall that $r
=
\frac{\Delta H(x)+\epsilon-v_0 T_\epsilon^\star(x)}
{L_H\bigl(1+e^{LT_\epsilon^\star(x)}\bigr)}$. Then
\[
\Delta H(x)+\epsilon-v_0 T_\epsilon(x)^\star
\geq
(\Delta H(x)+\epsilon)\left(1-\frac{v_0}{\underline{v_\epsilon}}\right)> 0,
\]
because of 
$0<v_0<\underline{v_\epsilon}$.

Therefore, for each fixed $x \in \X_c$, we obtain the lower bound of radius $r(x)$
\begin{align}
r(x) & 
\geq
\frac{
(\Delta H(x)+\epsilon)\left(1-\frac{v_0}{\underline{v_\epsilon}}\right)
}{
L_H\left(
1+\exp\!\left(\frac{L(\Delta H(x)+\epsilon)}{\underline{v_\epsilon}}\right)
\right)
} \notag \\
& \ge \frac{
\epsilon \left(1-\frac{v_0}{\underline{v_\epsilon}}\right)
}{2 L_H \exp(\frac{L(L_H D_{\mathcal{X}} +\epsilon)}{\underline{v_\epsilon}})
}, \label{lower_bound_r_i}
\end{align}
where equation~\eqref{lower_bound_r_i} holds from the fact that $H(x)$ is $L_H$-continuous.

\noindent\textbf{Step 3: Each certified ball covers a nontrivial energy interval.}
Fix $x\in\X_c$ and consider the ball $\cB_{r(x)}(x)$.
By Assumption~\ref{ass_strong_convexity}, for any $\|v\|=1$, we have
\begin{align*}
H(x+rv)+H(x-rv)\ge2H(x)+ \mu_H r^2,
\end{align*}
so at least one of $x\pm rv$ is larger than the average value, that is to say, 
\begin{align}
\label{eq:lower_bound}
\max \{H(x+rv),H(x-rv)\}\geq H(x)+\frac{\mu_H }{2}r^2.
\end{align}
For any $y_1, y_2$ and $y\in \cB_{r(x)}(x)$, according to~\eqref{eq:lower_bound} it follows that
\begin{align*}
& \max_{y_1,y_2\in \cB_{r(x)}(x)} \bigl(H(y_1)-H(y_2)\bigr)\\
& \ge \max_{y\in \cB_{r(x)}(x)} \bigl(H(y)-H(x)\bigr)\\
& \ge \frac{\mu_{H}}{2}r(x)^2.
\end{align*}
If $H_+^\star\notin H(\cB_{r(x)}(x))$, we have: 
\begin{align*}
    &\max_{y_1,y_2\in \cB_{r(x)}(x)} |\Delta H(y_1)-\Delta H(y_2)|\\
    &=\max_{y_1,y_2\in \cB_{r(x)}(x)} |H(y_1)-H(y_2)|\\
    &\geq \frac{\mu_{H}}{2}r(x)^2.
\end{align*}
And if $H_+^\star\in H(\cB_{r(x)}(x))$, since~\eqref{eq:lower_bound} holds, there exists at least one $y^\star \in \cB_{r(x)}(x)$ such that $H(x)+\frac{\mu_{H}}{2}r(x)^2 \le H(y^\star) \le \max H(\cB_{r(x)}(x))$. Thus $[H(x),H(x)+\frac{\mu_{H}}{2}r(x)^2]\subseteq H(\cB_{r(x)}(x))$, and therefore we have:
\begin{align*}
    \max\{|H(x)-H_+^\star|,|H(x)+\frac{\mu_{H}}{2}r(x)^2-H_+^\star|\}\geq \frac{\mu_{H}}{4}r(x)^2.
\end{align*}
Consequently, for any $y_1, y_2 \in \cB_{r(x)}(x)$, we have
\begin{align}
&\max_{y_1,y_2\in B_r(x)} |\Delta H(y_1)-\Delta H(y_2)| \notag\\
% & = \max_{y_1,y_2\in B_r(x)} |H(y_1) + H(y_2) - 2H_+^*|\\
& \ge \max_{y_1\in B_r(x)}|H(y_1)-H_+^\star| \label{eq_distance_lower_bound}\\
& \geq \frac{\mu_{H}}{4}r(x)^2,\notag
\end{align}
where inequality~\eqref{eq_distance_lower_bound} holds since we can chose $y_2 \in \cB_{r(x)}(x)$ such that $H(y_2) = H_+^*$. 

Above all, we know that $\forall x \in \X_c$, we have
\begin{align}
& \max_{y_1,y_2\in \cB_{r(x)}(x)} |\Delta H(y_1)- \Delta H(y_2)| \ge \frac{\mu_{H}}{4}r(x)^2 \notag\\
\ge & \frac{\mu_{H} (1-\frac{v_0}{\underline{v_\epsilon}})^2}{16L_H^2}\frac{\epsilon^2}{\exp(\frac{2L(L_HD_{\mathcal{X}}+\epsilon)}{\underline{v_\epsilon}})} \label{eq:energy_interval_lower}
\end{align}

\noindent\textbf{Step 4: Finite covering of the relevant energy interval.}
Let \[
[H_1,H_2]:=H(\X_c)\cup H(H_{\tgt}^\epsilon).
\]
For every $E\in[H_1,H_2]$, choose any $x\in\X_c$ with $H(x)=E$.
Then $E\in H(\cB_{r(x)}(x))$, so the family $\{H(\cB_{r(x)}(x))\}_{x\in\X_c}$ covers $[H_1,H_2]$.
Since $[H_1,H_2]$ is compact and 
\begin{align*}
& \max_{y_1,y_2\in \cB_{r(x)}(x)} |H(y_1)-H(y_2)| \geq \frac{\mu_{H}}{4}r(x)^2\\
\ge & \frac{\mu_{H} (1-\frac{v_0}{\underline{v_\epsilon}})^2}{16L_H^2}\frac{\epsilon^2}{\exp(\frac{2L(L_HD_{\mathcal{X}}+\epsilon)}{\underline{v_\epsilon}})}
\end{align*} there exists a finite subcover
\[
[H_1,H_2]\subseteq \bigcup_{i=1}^N \cB_{r_i}(x_i).
\]
For each selected center $x_i$, define $r_i:=r(x_i), u_i:=u_{x_i},$ and let
\[
\K:=\{(x_i,r_i,u_i)\}_{i=1}^N,
\]
where this finite selection implies
\[
[H_1,H_2]\subseteq H(\Supp(\K)),
\]
which verifies Condition~\ref{condition2} of
Theorem~\ref{thm_general_reachability}.

\noindent\textbf{Step 5: Ergodic Coverage.}
Moreover, under Assumption~\ref{ass_dim_ergodic_components},
each energy layer in the relevant range is itself a unique ergodic component.
Since $\Supp(\K)$ intersects every energy level in $[H_1,H_2]$,
Condition~\ref{condition3} of
Theorem~\ref{thm_general_reachability} also holds.

Therefore all three conditions of
Theorem~\ref{thm_general_reachability} are satisfied, and the resulting
chain policy $\pi_{\K}$ guarantees that for almost every $x_0\in S_0$,
there exists $t<\infty$ such that
\[
\phi(t,x_0,\pi_{\K})\in S_{\tgt}.
\]

\noindent\textbf{Step 6: Sample complexity bound.}
By~\eqref{eq:energy_interval_lower}, every selected interval $\cB_{r_i}(x_i)$
has length at least $\eta$. Hence a greedy interval-covering argument on
$[H_1,H_2]$ gives
\begin{align*}
& N \le \frac{H_2-H_1}{\frac{\mu_{H} (1-\frac{v_0}{\underline{v_\epsilon}})^2}{16L_H^2}\frac{\epsilon^2}{\exp(\frac{2L(L_HD_{\mathcal{X}}+\epsilon)}{\underline{v_\epsilon}})}}\\
& =
(H_2-H_1)\,
\frac{16L_H^2}
{\mu_H\left(1-\frac{v_0}{\underline v_\epsilon}\right)^2}
\frac{
\exp\!\left(
\frac{2L(L_HD_{\X}+\epsilon)}{\underline v_\epsilon}
\right)}
{\epsilon^2}.
\end{align*}
This completes the proof.
\end{proof}

\end{document}
